\section{Desarrollo Analítico}

En esta sección se analiza la propagación de los errores inherentes de los coeficientes $a$, $b$ y $c$ de forma analítica. Para ello, se considera a las raíces $x_1$ y $x_2$ como funciones de dichos coeficientes y se utiliza el cálculo diferencial para determinar la sensibilidad de los resultados ante variaciones en los datos de entrada.

\subsection{Derivadas Parciales de la Raíz $x_1$}

Se calculan las derivadas parciales de la función $x_1(a, b, c) = \frac{-b + \sqrt{b^2 - 4ac}}{2a}$ respecto a cada uno de sus coeficientes. Estas representan los coeficientes de sensibilidad:

\begin{itemize}
    \item Sensibilidad respecto al coeficiente $a$:
    $$\frac{\partial x_1}{\partial a} = \frac{2ca+b\sqrt{b^2-4ac}-b^2}{2a^2\sqrt{b^2-4ac}}$$

    \item Sensibilidad respecto al coeficiente $b$:
    $$\frac{\partial x_1}{\partial b} = \frac{-\sqrt{b^2-4ac}+b}{2a\sqrt{b^2-4ac}}$$

    \item Sensibilidad respecto al coeficiente $c$:
    $$\frac{\partial x_1}{\partial c} = \frac{-1}{\sqrt{b^2-4ac}}$$
\end{itemize}

\subsection{Propagación del Error Absoluto}

Utilizando el teorema de propagación de errores para el límite superior del error absoluto, se define $\Delta x_1$ como la suma de las contribuciones de cada error inherente ($i_a, i_b, i_c$) ponderadas por su respectiva sensibilidad:

$$\Delta x_1 \approx \left| \frac{\partial x_1}{\partial a} \right| \cdot i_a + \left| \frac{\partial x_1}{\partial b} \right| \cdot i_b + \left| \frac{\partial x_1}{\partial c} \right| \cdot i_c$$

Sustituyendo las expresiones obtenidas:

$$\Delta x_1 \approx \left| \frac{2ca+b\sqrt{b^2-4ac}-b^2}{2a^2\sqrt{b^2-4ac}} \right| \cdot i_a + \left| \frac{-\sqrt{b^2-4ac}+b}{2a\sqrt{b^2-4ac}} \right| \cdot i_b + \left| \frac{-1}{\sqrt{b^2-4ac}} \right| \cdot i_c$$

\subsection{Cálculo del Error Relativo}

El error relativo $e1_r$ se obtiene dividiendo el error absoluto estimado por el valor escalar de la propia raíz:

$$e1_r = \frac{\Delta x_1}{|x_1|}$$

Sustituyendo con la fórmula resolvente invertida para facilitar la simplificación:

$$e1_r \approx \left( \left| \frac{2ca+b\sqrt{b^2-4ac}-b^2}{2a^2\sqrt{b^2-4ac}} \right| \cdot i_a + \left| \frac{-\sqrt{b^2-4ac}+b}{2a\sqrt{b^2-4ac}} \right| \cdot i_b + \left| \frac{-1}{\sqrt{b^2-4ac}} \right| \cdot i_c \right) \cdot \left| \frac{2a}{-b + \sqrt{b^2 - 4ac}} \right|$$

\subsection{Análisis para la Raíz $x_2$}

De manera análoga, para la segunda raíz $x_2(a, b, c) = \frac{-b - \sqrt{b^2 - 4ac}}{2a}$, se obtienen las expresiones de sensibilidad y error relativo correspondientes:

$$e2_r \approx \left( \left| \frac{-2ca-b\sqrt{b^2-4ac}+b^2}{2a^2\sqrt{b^2-4ac}} \right| \cdot i_a + \left| \frac{-\sqrt{b^2-4ac}-b}{2a\sqrt{b^2-4ac}} \right| \cdot i_b + \left| \frac{1}{\sqrt{b^2-4ac}} \right| \cdot i_c \right) \cdot \left| \frac{2a}{-b - \sqrt{b^2 - 4ac}} \right|$$

